% LTeX: language=de

\subsection{Ober- und Untersumme}

\begin{Example}{Fläche unter dem Graphen}{}
    \begin{equation*}
        f(x) = x^{2}
    \end{equation*}
    Bestimme die Fläche zwischen $G_{f}$ und der x-Achse im Intervall $I = [ \ 0; 2 \ ]$.
\end{Example}

\begin{solution}{}
    \begin{itemize}[leftmargin=*]
        \item in 12: \quad {
            \begin{minipage}[t]{0.7\linewidth}
                \vspace{-2.0ex}
                \begin{align*}
                    A &= \int_{0}^{2} f(x) \,\mathrm{d}x = \bigl[ \ F(x) \ \bigr]^{2}_{0}\\[1.5ex]
                    &= \left[ \ \dfrac{1}{3} x^{3} \ \right]^{2}_{0} = \dfrac{1}{3} \cdot 2^{3} - \dfrac{1}{3} \cdot 0^{3} = \dfrac{8}{3}
                \end{align*}
            \end{minipage}
            }
        \item in 13: \quad {
            \begin{minipage}[t]{0.85\linewidth}
                \begin{minipage}[t]{0.45\linewidth}
                    \centering
                    \raisebox{-\height}{
                    \begin{tikzpicture}[
                            declare function={
                                f_f(\x) = ((\x*\xScale)^2)*(1/\yScale);
                            }
                        ]

                        % Set variables
                        \pgfmathsetmacro{\axisPosWidth}{4.5};
                        \pgfmathsetmacro{\axisNegWidth}{0.0};
                        \pgfmathsetmacro{\axisPosHeight}{5.0};
                        \pgfmathsetmacro{\axisNegHeight}{0.0};
                        \pgfmathsetmacro{\xIncrement}{2.0};
                        \pgfmathsetmacro{\yIncrement}{1.0};
                        \pgfmathsetmacro{\xScale}{0.5};
                        \pgfmathsetmacro{\yScale}{1.0};
                        \pgfmathsetmacro{\gridxIncrement}{0.5};
                        \pgfmathsetmacro{\gridyIncrement}{0.5};

                        % Axis/grid limits
                        \pgfmathsetmacro{\xUpperLimitAxis}{\xIncrement * floor(\axisPosWidth / \xIncrement)}
                        \pgfmathsetmacro{\xLowerLimitAxis}{\xIncrement * ceil(\axisNegWidth / \xIncrement)}
                        \pgfmathsetmacro{\yUpperLimitAxis}{\yIncrement * floor(\axisPosHeight /\yIncrement)}
                        \pgfmathsetmacro{\yLowerLimitAxis}{\yIncrement * ceil(\axisNegHeight / \yIncrement)}

                        \pgfmathsetmacro{\xUpperLimitGrid}{\gridxIncrement * floor(\axisPosWidth / \gridxIncrement)}
                        \pgfmathsetmacro{\xLowerLimitGrid}{\gridxIncrement * ceil(\axisNegWidth / \gridxIncrement)}
                        \pgfmathsetmacro{\yUpperLimitGrid}{\gridyIncrement * floor(\axisPosHeight / \gridyIncrement)}
                        \pgfmathsetmacro{\yLowerLimitGrid}{\gridyIncrement * ceil(\axisNegHeight / \gridyIncrement)}

                        % Axis environment (PGFPlots)
                        \begin{axis}[
                            xmin=\xLowerLimitGrid, xmax=\xUpperLimitGrid,
                            ymin=\yLowerLimitGrid, ymax=\yUpperLimitGrid,
                            axis lines=middle,
                            axis line style={->, >=Stealth},
                            grid=both,
                            xlabel={$x$},
                            ylabel={$y$},
                            xlabel style={anchor=west, at={(current axis.right of origin)}},
                            ylabel style={anchor=south, at={(current axis.above origin)}},
                            xtick={\xLowerLimitAxis,\xLowerLimitAxis+1,...,\xUpperLimitAxis},
                            ytick={\yLowerLimitAxis,\yLowerLimitAxis+1,...,\yUpperLimitAxis},
                            xticklabel={\pgfmathparse{\tick*\xScale}\pgfmathprintnumber{\pgfmathresult}},
                            yticklabel={\pgfmathparse{\tick*\yScale}\pgfmathprintnumber{\pgfmathresult}},
                            tick label style={font=\scriptsize},
                            x=1cm,
                            y=1cm,
                            minor tick num=1,
                            scale only axis,
                            grid=both,
                            restrict y to domain=\yLowerLimitGrid:\yUpperLimitGrid,
                        ]
                            
                        % Graph limits
                        \pgfmathsetmacro{\xLowerLimitGraph}{\xLowerLimitAxis}
                        \pgfmathsetmacro{\xUpperLimitGraph}{\xUpperLimitAxis}
                        
                        % Interval and subdivisions
                        \pgfmathsetmacro{\a}{0*(1/\xScale)}
                        \pgfmathsetmacro{\b}{2*(1/\xScale)}
                        \pgfmathsetmacro{\n}{4}
                        \pgfmathsetmacro{\dx}{(\b - \a)/\n}
                        \pgfmathsetmacro{\nMinusOne}{\n-1}

                        % G_f
                        \pgfmathsetmacro{\Gfax}{2*(1/\xScale)}
                        \pgfmathsetmacro{\Gfay}{f_f(\Gfax)}
                        \addplot[thick, myb, samples=1000, domain=\xLowerLimitGraph:\xUpperLimitGraph]
                            {f_f(x)};
                        \node[anchor=south] at (axis cs:{\Gfax},{\Gfay}) {\textcolor{myb}{$G_{f}$}};

                        % ----- UPPER SUM -----
                        \foreach \i in {0,...,\nMinusOne} {
                            \pgfmathsetmacro{\xleft}{\a + \i*\dx}
                            \pgfmathsetmacro{\xright}{\xleft + \dx}
                            \pgfmathsetmacro{\yleft}{f_f(\xleft)}
                            \pgfmathsetmacro{\yright}{f_f(\xright)}
                            \pgfmathsetmacro{\ymax}{max(\yleft,\yright)}

                            \addplot[fill=mylightgreen, draw=myg, opacity=0.6]
                            coordinates {
                                (\xleft,0) (\xleft,\ymax) (\xright,\ymax) (\xright,0)
                            } -- cycle;
                        }

                        \end{axis}
                    \end{tikzpicture}
                    }
                    \textcolor{myg}{\textbf{Obersumme}}
                \end{minipage}
                \begin{minipage}[t]{0.40\linewidth}
                    \centering
                    \raisebox{-\height}{
                    \begin{tikzpicture}[
                            declare function={
                                f_f(\x) = ((\x*\xScale)^2)*(1/\yScale);
                            }
                        ]

                        % Set variables
                        \pgfmathsetmacro{\axisPosWidth}{4.5};
                        \pgfmathsetmacro{\axisNegWidth}{0.0};
                        \pgfmathsetmacro{\axisPosHeight}{5.0};
                        \pgfmathsetmacro{\axisNegHeight}{0.0};
                        \pgfmathsetmacro{\xIncrement}{2.0};
                        \pgfmathsetmacro{\yIncrement}{1.0};
                        \pgfmathsetmacro{\xScale}{0.5};
                        \pgfmathsetmacro{\yScale}{1.0};
                        \pgfmathsetmacro{\gridxIncrement}{0.5};
                        \pgfmathsetmacro{\gridyIncrement}{0.5};

                        % Axis/grid limits
                        \pgfmathsetmacro{\xUpperLimitAxis}{\xIncrement * floor(\axisPosWidth / \xIncrement)}
                        \pgfmathsetmacro{\xLowerLimitAxis}{\xIncrement * ceil(\axisNegWidth / \xIncrement)}
                        \pgfmathsetmacro{\yUpperLimitAxis}{\yIncrement * floor(\axisPosHeight /\yIncrement)}
                        \pgfmathsetmacro{\yLowerLimitAxis}{\yIncrement * ceil(\axisNegHeight / \yIncrement)}

                        \pgfmathsetmacro{\xUpperLimitGrid}{\gridxIncrement * floor(\axisPosWidth / \gridxIncrement)}
                        \pgfmathsetmacro{\xLowerLimitGrid}{\gridxIncrement * ceil(\axisNegWidth / \gridxIncrement)}
                        \pgfmathsetmacro{\yUpperLimitGrid}{\gridyIncrement * floor(\axisPosHeight / \gridyIncrement)}
                        \pgfmathsetmacro{\yLowerLimitGrid}{\gridyIncrement * ceil(\axisNegHeight / \gridyIncrement)}

                        % Axis environment (PGFPlots)
                        \begin{axis}[
                            xmin=\xLowerLimitGrid, xmax=\xUpperLimitGrid,
                            ymin=\yLowerLimitGrid, ymax=\yUpperLimitGrid,
                            axis lines=middle,
                            axis line style={->, >=stealth},
                            grid=both,
                            xlabel={$x$},
                            ylabel={$y$},
                            xlabel style={anchor=west, at={(current axis.right of origin)}},
                            ylabel style={anchor=south, at={(current axis.above origin)}},
                            xtick={\xLowerLimitAxis,\xLowerLimitAxis+1,...,\xUpperLimitAxis},
                            ytick={\yLowerLimitAxis,\yLowerLimitAxis+1,...,\yUpperLimitAxis},
                            xticklabel={\pgfmathparse{\tick*\xScale}\pgfmathprintnumber{\pgfmathresult}},
                            yticklabel={\pgfmathparse{\tick*\yScale}\pgfmathprintnumber{\pgfmathresult}},
                            tick label style={font=\scriptsize},
                            x=1cm,
                            y=1cm,
                            minor tick num=1,
                            scale only axis,
                            grid=both,
                            restrict y to domain=\yLowerLimitGrid:\yUpperLimitGrid,
                        ]
                            
                        % Graph limits
                        \pgfmathsetmacro{\xLowerLimitGraph}{\xLowerLimitAxis}
                        \pgfmathsetmacro{\xUpperLimitGraph}{\xUpperLimitAxis}
                        
                        % Interval and subdivisions
                        \pgfmathsetmacro{\a}{0*(1/\xScale)}
                        \pgfmathsetmacro{\b}{2*(1/\xScale)}
                        \pgfmathsetmacro{\n}{4}
                        \pgfmathsetmacro{\dx}{(\b - \a)/\n}
                        \pgfmathsetmacro{\nMinusOne}{\n-1}

                        % G_f
                        \pgfmathsetmacro{\Gfax}{2*(1/\xScale)}
                        \pgfmathsetmacro{\Gfay}{f_f(\Gfax)}
                        \addplot[thick, myb, samples=1000, domain=\xLowerLimitGraph:\xUpperLimitGraph]
                            {f_f(x)};
                        \node[anchor=south] at (axis cs:{\Gfax},{\Gfay}) {\textcolor{myb}{$G_{f}$}};

                        % ----- LOWER SUM -----
                        \foreach \i in {0,...,\nMinusOne} {
                            \pgfmathsetmacro{\xleft}{\a + \i*\dx}
                            \pgfmathsetmacro{\xright}{\xleft + \dx}
                            \pgfmathsetmacro{\yleft}{f_f(\xleft)}
                            \pgfmathsetmacro{\yright}{f_f(\xright)}
                            \pgfmathsetmacro{\ymin}{min(\yleft,\yright)}

                            \addplot[fill=mylightorange, draw=myorange, opacity=0.6]
                            coordinates {
                                (\xleft,0) (\xleft,\ymin) (\xright,\ymin) (\xright,0)
                            } -- cycle;
                        }

                        \end{axis}
                    \end{tikzpicture}
                    }
                    \textcolor{myorange}{\textbf{Untersumme}}
                \end{minipage}\\[1.0ex]
                \begin{enumerate}[leftmargin=*, label=\Roman*)]
                    \item Unterteile $I$ in gleich große Teilintervalle (hier z. B. $n = 4$) der Breite $\Delta x$ (hier $\Delta x = \SI{0.5}{}$)
                    \item Überdecke die gesuchte Fläche mit Rechtecken so dass \\[-2.0ex]{
                        \begin{itemize}[leftmargin=*]
                            \item \textcolor{myg}{die Fläche vollständig abgedeckt ist (\textbf{Obersumme})}
                            \item \textcolor{myorange}{jedes Rechteck vollständig innerhalb der gesuchten Fläche liegt (\textbf{Untersumme})}
                        \end{itemize}
                        }
                    \item $\textcolor{myorange}{\text{Untersumme} \ U} \quad \leq A \quad \leq \textcolor{myg}{\text{Obersumme} \ O}$
                    \item Je kleiner $\Delta x$ (d. h. je mehr Rechtecke), umso genauer wird die Abschätzung
                \end{enumerate}
            \end{minipage}
            }
    \end{itemize}

\end{solution}
